\section{模型假设}

\begin{enumerate}
\item 药材近似为均匀、各向同性的长圆柱，长度 $L=0.25\,\mathrm{m}$；模型只描述径向传热传质，忽略周向和轴向变化，主要适用于圆柱主体区域。
\item 问题一至三半径保持不变；问题四长度不变，仅发生径向收缩，材料位置随半径同比例缩放，并用材料坐标描述几何运动。
\item 初始温度和干基水分浓度在药材内部均匀分布，温度采用 Fourier 导热方程，水分采用 Fick 扩散方程；表面采用第三类对流边界，轴线采用轴对称零通量条件。
\item 附件1环境数据在给定时间范围内分段线性插值；长时段计算超出数据范围时，环境量保持为附件末值。问题四的半径轨迹由附件2数据插值确定。
\item 本文不显式考虑蒸发潜热、吸附等温线、Soret/Dufour 效应及未经题面给出的收缩本构关系。这些是模型简化，不表示相应机制在实际过程中不存在。
\end{enumerate}

上述假设限定了模型的适用范围；网格、时间积分和表面重构属于数值实现，在各问题的求解与检验部分说明。
